\chapter{Numbers and Sets of Numbers}\label{ch:g10:numbers}

Mathematics begins with numbers, and not all numbers are of the same kind:
counting numbers, negative numbers, fractions, and numbers like $\sqrt 2$
or $\pi$ that no fraction can express. This chapter organizes them into
nested families, introduces \omterm{def:g10:numbers:interval}{intervals} to describe portions of the number
line, and uses the \omterm{def:g10:numbers:abs}{absolute value} to measure distances between numbers.

\section{The families of numbers}

\begin{definition}[Sets of numbers]\label{def:g10:numbers:sets}
\begin{itemize}
  \item $\N$ is the set of \emph{natural numbers}\index{natural number}:
        $0, 1, 2, 3, \dots$
  \item $\Z$ is the set of \emph{integers}\index{integer}: $\dots, -2, -1,
        0, 1, 2, \dots$
  \item $\Q$ is the set of \emph{rational numbers}\index{rational number}:
        all quotients $\frac{p}{q}$ with $p \in \Z$, $q \in \N$ and
        $q \neq 0$.
  \item $\R$ is the set of \emph{real numbers}\index{real number}: all the
        numbers that can be placed on the number line.
\end{itemize}
\end{definition}

\begin{notation}
The symbol $\in$ reads ``belongs to'': $3 \in \N$, $-5 \in \Z$,
$\frac23 \in \Q$. The symbol $\subset$ reads ``is included in'': every
\omterm{def:g10:numbers:sets}{natural number} is an \omterm{def:g10:numbers:sets}{integer}, every \omterm{def:g10:numbers:sets}{integer} is a \omterm{def:g10:numbers:sets}{rational number}
(e.g.\ $-5 = \frac{-5}{1}$), and every \omterm{def:g10:numbers:sets}{rational number} is real, so
\[
\N \subset \Z \subset \Q \subset \R .
\]
A slash negates a symbol: $\frac12 \notin \Z$.
\end{notation}

\begin{omfigure}
\begin{tikzpicture}[scale=0.9]
  \draw[omExo, thick] (0,0) ellipse (5.4 and 2.9);
  \draw[omProp, thick] (-0.8,0) ellipse (4.0 and 2.2);
  \draw[omThm, thick] (-1.6,0) ellipse (2.6 and 1.5);
  \draw[omDef, thick] (-2.4,0) ellipse (1.3 and 0.8);
  \node[omDef, font=\small\bfseries] at (-2.4,0.45) {$\N$};
  \node[omThm, font=\small\bfseries] at (-1.6,1.1) {$\Z$};
  \node[omProp, font=\small\bfseries] at (-0.8,1.8) {$\Q$};
  \node[omExo, font=\small\bfseries] at (0,2.5) {$\R$};
  \node[font=\small] at (-2.4,-0.15) {$0 \quad 7$};
  \node[font=\small] at (-0.4,-0.9) {$-5$};
  \node[font=\small] at (1.0,-0.4) {$\tfrac23 \quad {-1.7}$};
  \node[font=\small] at (3.9,-0.6) {$\sqrt2 \quad \pi$};
\end{tikzpicture}

{\small The nested families of numbers: $\N \subset \Z \subset \Q \subset
\R$. Each ring contains numbers that do not belong to the smaller ones.}
\end{omfigure}

\begin{example}\label{ex:g10:numbers:classify}
Let us place a few numbers in the smallest family that contains them.
\begin{itemize}
  \item $\frac{28}{4} = 7$, so $\frac{28}{4} \in \N$ even though it is
        written as a fraction: always simplify first.
  \item $-3.5 = -\frac{35}{10} = -\frac72$ is rational but not an \omterm{def:g10:numbers:sets}{integer}.
  \item $0.333\dots$ (the digit $3$ repeating forever) equals $\frac13$,
        a \omterm{def:g10:numbers:sets}{rational number}.
  \item $\sqrt 2$ and $\pi$ are real but not rational, as we are about to
        see for $\sqrt 2$; such numbers are called
        \emph{irrational}\index{irrational number}.
\end{itemize}
\end{example}

\begin{theorem}[Irrationality of $\sqrt 2$]\label{thm:g10:numbers:sqrt2}
The number $\sqrt 2$ is not rational: no fraction of \omterm{def:g10:numbers:sets}{integers} has square
$2$.
\end{theorem}

\begin{proof}
We reason by contradiction, in small steps.
\begin{enumerate}
  \item Suppose $\sqrt 2 = \frac{p}{q}$ where $p$ and $q$ are positive
        \omterm{def:g10:numbers:sets}{integers}, and the fraction is fully simplified, so $p$ and $q$ are
        not both even.
  \item Squaring both sides gives $2 = \frac{p^2}{q^2}$, that is
        $p^2 = 2q^2$. So $p^2$ is even.
  \item If $p$ were odd, say $p = 2k+1$, then
        $p^2 = 4k^2 + 4k + 1$ would be odd. Since $p^2$ is even, $p$ must
        be even: $p = 2k$ for some \omterm{def:g10:numbers:sets}{integer} $k$.
  \item Substituting: $(2k)^2 = 2q^2$, so $4k^2 = 2q^2$, so
        $q^2 = 2k^2$. By the same argument as in step 3, $q$ is even.
  \item Now $p$ and $q$ are both even, contradicting step 1. The
        assumption was impossible: $\sqrt 2$ is \omterm{ex:g10:numbers:classify}{irrational}.
\end{enumerate}
\end{proof}

\begin{remark}
\omterm{def:g10:numbers:sets}{Rational numbers} are exactly the numbers whose decimal expansion either
stops (like $\frac72 = 3.5$) or eventually repeats the same block forever
(like $\frac13 = 0.333\dots$ or $\frac{1}{7} = 0.142857\,142857\dots$).
\omterm{ex:g10:numbers:classify}{Irrational numbers} like $\sqrt2 = 1.41421356\dots$ have decimal digits
that never fall into a repeating pattern. We admit this characterization
at this level.
\end{remark}

\section{The number line and intervals}

\omterm{def:g10:numbers:sets}{Real numbers} fill a line: choosing an origin $0$ and a unit of length,
every \omterm{def:g10:numbers:sets}{real number} corresponds to exactly one point.

\begin{definition}[Interval]\label{def:g10:numbers:interval}
Let $a$ and $b$ be \omterm{def:g10:numbers:sets}{real numbers} with $a < b$. An
\emph{interval}\index{interval} is the set of all \omterm{def:g10:numbers:sets}{real numbers} between two
bounds. A square bracket means the bound is included, a parenthesis means
it is excluded:
\begin{center}
\begin{tabular}{c|c}
notation & description \\
\hline
$\intcc{a}{b}$ & $a \leq x \leq b$ (both ends included) \\
$\intoo{a}{b}$ & $a < x < b$ (both ends excluded) \\
$\intco{a}{b}$ & $a \leq x < b$ \\
$\intoc{a}{b}$ & $a < x \leq b$ \\
$\intco{a}{+\infty}$ & $x \geq a$ \\
$\intoo{-\infty}{b}$ & $x < b$ \\
\end{tabular}
\end{center}
The symbols $-\infty$ and $+\infty$ (``infinity'') are not numbers, only a
way of saying that the interval continues forever; the bracket next to
them is always a parenthesis. The whole line $\R$ is the interval
$\intoo{-\infty}{+\infty}$.
\end{definition}

\begin{omfigure}
\begin{tikzpicture}[scale=1.0]
  \draw[-Stealth] (-4.4,0) -- (4.6,0) node[below, font=\small] {$x$};
  \foreach \x in {-4,...,4} \draw (\x,-0.07) -- (\x,0.07);
  \foreach \x in {-4,-2,0,2,4}
    \node[below, font=\small] at (\x,-0.1) {$\x$};
  \draw[omDef, very thick] (-3,0.02) -- (1,0.02);
  \fill[omDef] (-3,0.02) circle (2.2pt);
  \draw[omDef, fill=white, thick] (1,0.02) circle (2.2pt);
  \node[omDef, above, font=\small] at (-1,0.15) {$\intco{-3}{1}$};
  \draw[omThm, very thick] (2,0.02) -- (4.5,0.02);
  \draw[omThm, fill=white, thick] (2,0.02) circle (2.2pt);
  \node[omThm, above, font=\small] at (3.4,0.15) {$\intoo{2}{+\infty}$};
\end{tikzpicture}

{\small \omterm{def:g10:numbers:interval}{Intervals} on the number line: a filled dot for an included bound,
a hollow dot for an excluded one.}
\end{omfigure}

\begin{definition}[Intersection and union]\label{def:g10:numbers:interunion}
Let $I$ and $J$ be two sets of \omterm{def:g10:numbers:sets}{real numbers}.
The \emph{intersection}\index{intersection} $I \cap J$ (``$I$ and $J$'')
is the set of numbers belonging to both; the
\emph{union}\index{union} $I \cup J$ (``$I$ or $J$'') is the set of
numbers belonging to at least one of them.
\end{definition}

\begin{example}\label{ex:g10:numbers:interunion}
Take $I = \intcc{-1}{3}$ and $J = \intoo{1}{5}$. Draw both on the same
line: they overlap between $1$ and $3$. Therefore
\[
I \cap J = \intoc{1}{3},
\qquad
I \cup J = \intco{-1}{5}.
\]
Note the brackets: $1 \notin J$ so $1 \notin I \cap J$, but $3 \in I$ and
$3 \in J$, so $3 \in I \cap J$.
\end{example}

\begin{method}[Working with intervals]\label{met:g10:numbers:intervals}
To find the \omterm{def:g10:numbers:interunion}{intersection} or \omterm{def:g10:numbers:interunion}{union} of two \omterm{def:g10:numbers:interval}{intervals}:
\begin{enumerate}
  \item draw the number line and mark the four bounds;
  \item shade the first \omterm{def:g10:numbers:interval}{interval} above the line and the second below it;
  \item the \omterm{def:g10:numbers:interunion}{intersection} is where the shadings overlap, the \omterm{def:g10:numbers:interunion}{union} is where
        at least one shading is present;
  \item decide each bracket by testing whether the bound itself belongs to
        both sets (\omterm{def:g10:numbers:interunion}{intersection}) or to at least one (\omterm{def:g10:numbers:interunion}{union}).
\end{enumerate}
\end{method}

\section{Absolute value and distance}

\begin{definition}[Absolute value]\label{def:g10:numbers:abs}
The \emph{absolute value}\index{absolute value} of a \omterm{def:g10:numbers:sets}{real number} $x$ is
\[
\abs{x} =
\begin{cases}
x & \text{if } x \geq 0, \\
-x & \text{if } x < 0.
\end{cases}
\]
For instance $\abs{7} = 7$ and $\abs{-4} = 4$. In every case
$\abs{x} \geq 0$.
\end{definition}

\begin{proposition}[Distance on the line]\label{prop:g10:numbers:distance}
For all \omterm{def:g10:numbers:sets}{real numbers} $a$ and $b$, the distance between the points $a$ and
$b$ on the number line is $\abs{b - a}$. In particular $\abs{x}$ is the
distance from $x$ to $0$.
\end{proposition}

\begin{proof}
If $b \geq a$, the distance from $a$ to $b$ is $b - a \geq 0$, which
equals $\abs{b-a}$. If $b < a$, the distance is $a - b = -(b - a) > 0$,
which is again $\abs{b - a}$ by definition of the \omterm{def:g10:numbers:abs}{absolute value}.
\end{proof}

\begin{proposition}[Absolute value and intervals]\label{prop:g10:numbers:absinterval}
Let $a$ be a \omterm{def:g10:numbers:sets}{real number} and $r > 0$. Then
\[
\abs{x - a} \leq r
\quad\text{exactly when}\quad
x \in \intcc{a - r}{a + r}.
\]
\end{proposition}

\begin{proof}
$\abs{x-a} \leq r$ says that the distance from $x$ to $a$ is at most $r$,
i.e.\ that $x$ lies no further than $r$ from $a$ on either side. The
numbers satisfying this are exactly those between $a - r$ and $a + r$,
bounds included.
\end{proof}

\begin{omfigure}
\begin{tikzpicture}[scale=1.1]
  \draw[-Stealth] (-0.8,0) -- (7,0) node[below, font=\small] {$x$};
  \draw[omDef, very thick] (1.5,0.02) -- (5.5,0.02);
  \fill[omDef] (1.5,0.02) circle (2.2pt);
  \fill[omDef] (5.5,0.02) circle (2.2pt);
  \fill (3.5,0) circle (1.6pt);
  \node[below, font=\small] at (3.5,-0.08) {$a$};
  \node[below, font=\small] at (1.5,-0.08) {$a-r$};
  \node[below, font=\small] at (5.5,-0.08) {$a+r$};
  \draw[Stealth-Stealth, omThm] (3.5,0.35) -- (5.5,0.35)
    node[midway, above, font=\small] {$r$};
  \draw[Stealth-Stealth, omThm] (1.5,0.35) -- (3.5,0.35)
    node[midway, above, font=\small] {$r$};
\end{tikzpicture}

{\small The inequality $\abs{x - a} \leq r$ describes the \omterm{def:g10:numbers:interval}{interval} of
numbers at distance at most $r$ from $a$.}
\end{omfigure}

\begin{example}\label{ex:g10:numbers:abs}
Solve $\abs{x - 3} \leq 2$. The solutions are the numbers at distance at
most $2$ from $3$: the \omterm{def:g10:numbers:interval}{interval} $\intcc{1}{5}$. Conversely, the \omterm{def:g10:numbers:interval}{interval}
$\intcc{-1}{7}$ has center $\frac{-1+7}{2} = 3$ and radius
$\frac{7-(-1)}{2} = 4$, so it is described by $\abs{x - 3} \leq 4$.
\end{example}

\section{Approximations}

\omterm{ex:g10:numbers:classify}{Irrational numbers}, and even most fractions, cannot be written exactly
with finitely many decimal digits, so in practice we approximate them.

\begin{definition}[Approximation to a given accuracy]\label{def:g10:numbers:approx}
A number $d$ is an \emph{approximation}\index{approximation} of $x$ to
within $10^{-n}$ when $\abs{x - d} \leq 10^{-n}$. Truncating or rounding
the decimal expansion after the $n$-th digit both give such
approximations.
\end{definition}

\begin{example}
From $\pi = 3.14159\,26\dots$: the truncation $3.141$ and the rounding
$3.142$ are both \omterm{def:g10:numbers:approx}{approximations} of $\pi$ to within $10^{-3}$. The rounding
is at distance less than $\frac12 \times 10^{-3}$, the truncation only
guarantees $10^{-3}$. Writing $3.141 \leq \pi \leq 3.142$ \emph{frames}
$\pi$ between two decimal bounds.
\end{example}

\section{Exercises}

\begin{exercise}[$\star$]\label{exo:g10:numbers:1}
For each number, give the smallest of the sets $\N$, $\Z$, $\Q$, $\R$ to
which it belongs:
\[
\frac{15}{3}, \qquad -7, \qquad \frac{22}{7}, \qquad \sqrt{9}, \qquad
\sqrt{10}, \qquad -2.4, \qquad 0 .
\]
\end{exercise}

\begin{exercise}[$\star$]\label{exo:g10:numbers:2}
Write each statement with \omterm{def:g10:numbers:interval}{interval} notation, then draw it on a number
line: (a) $-2 \leq x < 5$; (b) $x > 3$; (c) $x \leq -1$;
(d) the distance from $x$ to $2$ is at most $3$.
\end{exercise}

\begin{exercise}[$\star$]\label{exo:g10:numbers:3}
Compute $I \cap J$ and $I \cup J$ for
\[
\text{(a) } I = \intcc{-3}{2},\ J = \intco{0}{4};
\qquad
\text{(b) } I = \intoo{-\infty}{1},\ J = \intco{-2}{+\infty}.
\]
\end{exercise}

\begin{exercise}[$\star$]\label{exo:g10:numbers:4}
Compute without a calculator:
\[
\abs{-6}, \qquad \abs{4 - 9}, \qquad \abs{-3 - 5}, \qquad
\abs{\sqrt 2 - 1}, \qquad \abs{1 - \sqrt 2}.
\]
\end{exercise}

\begin{exercise}[$\star$]\label{exo:g10:numbers:5}
Solve the equations and inequalities and give the solution sets:
\[
\abs{x} = 5, \qquad \abs{x - 1} = 3, \qquad \abs{x - 4} \leq 1, \qquad
\abs{x + 2} < 3 .
\]
(Note that $\abs{x+2} = \abs{x - (-2)}$ is a distance to $-2$.)
\end{exercise}

\begin{exercise}[$\star\star$]\label{exo:g10:numbers:6}
Describe each \omterm{def:g10:numbers:interval}{interval} by an inequality of the form
$\abs{x - a} \leq r$ or $\abs{x - a} < r$:
\[
\intcc{2}{8}, \qquad \intoo{-5}{1}, \qquad \intcc{-7}{-3}.
\]
\end{exercise}

\begin{exercise}[$\star\star$]\label{exo:g10:numbers:7}
Show that $0.272727\dots$ (the block $27$ repeating forever) is rational.
(Hint: call it $x$ and compute $100x - x$.)
\end{exercise}

\begin{exercise}[$\star\star$]\label{exo:g10:numbers:8}
True or false? Justify each answer with an argument or a counterexample.
\begin{enumerate}
  \item The sum of two \omterm{def:g10:numbers:sets}{integers} is an \omterm{def:g10:numbers:sets}{integer}.
  \item The quotient of two \omterm{def:g10:numbers:sets}{integers} is an \omterm{def:g10:numbers:sets}{integer}.
  \item The sum of two \omterm{def:g10:numbers:sets}{rational numbers} is rational.
  \item The sum of a \omterm{def:g10:numbers:sets}{rational number} and an \omterm{ex:g10:numbers:classify}{irrational number} is
        \omterm{ex:g10:numbers:classify}{irrational}.
\end{enumerate}
\end{exercise}

\begin{exercise}[$\star\star$]\label{exo:g10:numbers:9}
Using $1.414 \leq \sqrt 2 \leq 1.415$, frame the numbers $2\sqrt2$,
$\sqrt2 + 3$ and $-\sqrt2$ between two decimal bounds.
\end{exercise}

\begin{exercise}[$\star\star\star$]\label{exo:g10:numbers:10}
Adapt the proof of \cref{thm:g10:numbers:sqrt2} to show that $\sqrt 3$ is
\omterm{ex:g10:numbers:classify}{irrational}. (Replace ``even'' by ``multiple of $3$'': first check that if
$p^2$ is a multiple of $3$, then so is $p$, by examining the remainders
$0$, $1$, $2$ of $p$ upon division by $3$.)
\end{exercise}
